% Separate refinements: the original height-two candidate remains byte-stable.
\subsection*{What the height-two countermodel does and does not obstruct}

Use \(N=2m+1\), \(m\ge2\), and the tensor \(\beta\) and reduced
radical \(\sigma\) of (PH2.1)--(PH2.7).

\begin{lemma}[The compatibility polynomial is nevertheless irreducible]
\label{lem:abstract-height-two-q-irreducible}
\status{Candidate; direct Artin--Schreier proof.}
Over an algebraically closed field of characteristic two, the polynomial
\[
                    q(x)=\sum_{i=0}^{2m}x_i\sigma_i(x)^2
\]
is irreducible for this countermodel.
\end{lemma}
\begin{proof}
On the chart \(x_1=1\), write \(u=x_0\), \(z=x_2\), and
\[
 c=x_{2m-1}x_{2m}(x_{2m-1}+x_{2m}),\qquad
 d=\sum_{j=2}^{m-1}x_{2j-1}x_{2j}(x_{2j-1}+x_{2j}).
\]
The empty sum for \(m=2\) is zero.  Formula (PH2.3) gives
\[
                        q=z^2+z+u^3+cu^2+d.                    \tag{PH2R.1}
\]
Let \(E\) be the rational function field in all the other coordinates.
At infinity in \(E(u)\), the element \(u^3+cu^2+d\) has pole
order three.  For any \(b\in E(u)\), either \(b\) has no pole
there, or the pole order of \(b^2+b\) is twice the pole order of
\(b\).  Hence \(u^3+cu^2+d\) is not \(b^2+b\).
The quadratic polynomial (PH2R.1) is therefore irreducible over
\(E(u)\), and Gauss's lemma makes it irreducible in the affine
coordinate ring.  The original homogeneous polynomial has no
\(x_1\) factor, because it contains
\(x_0^2x_{2m-1}x_{2m}^2\).  An irreducible dehomogenization
with no such factor proves homogeneous irreducibility.
\end{proof}

\begin{lemma}[Its linear centroid is scalar]
\label{lem:abstract-height-two-scalar-centroid}
\status{Candidate; direct coefficient comparison.}
If \(D\in\operatorname{End}(k^N)\) satisfies
\(\beta(s,Ds)=0\) identically as a polynomial in \(s\), then
\(D\) is a scalar.
\end{lemma}
\begin{proof}
For every \(i\ge2\), the \(i\)-th output equation in (PH2.7) is
\[
                       s_0(Ds)_i+s_i(Ds)_0=0.                 \tag{PH2R.2}
\]
The coefficient of \(s_i^2\) gives \(D_{0i}=0\); the coefficient
of \(s_is_1\) gives \(D_{01}=0\).  The coefficient of \(s_0^2\)
gives \(D_{i0}=0\), and the remaining mixed coefficients give
\(D_{ij}=0\) for \(j\ne i\) and \(D_{ii}=D_{00}\).
Thus row zero and every row \(i\ge2\) are those of \(aI\),
where \(a=D_{00}\).
The zeroth output equation now reduces to
\[
                         s_2\bigl((Ds)_1+a s_1\bigr)=0.
\]
The polynomial ring is a domain, so row one is also the row of
\(aI\).  Conversely scalar \(D\) clearly satisfies the condition
by alternation.
\end{proof}

\paragraph{A genuine missing equivariance condition.}
Suppose a group acts linearly and irreducibly on the parameter
space \(k^N\), and makes the pencil equivariant up to an invertible
congruence and a nonzero scalar.  The intrinsic radical rational
map then has an invariant reduced indeterminacy support.
In the present example that support has exactly one irreducible
component of codimension two, namely \(V(x_0,x_1)\).
Every group element must preserve that component, because it
cannot permute it with the codimension-three component.
Its corresponding nonzero proper vector subspace would be
invariant, contradicting irreducibility.  Thus this example
cannot carry such an irreducible parameter action.

The actual odd theta-group action is an additional geometric
constraint of this kind; the countermodel is not claimed to
realize it.  This does not prove that theta equivariance alone
forces every possible reduced Pfaffian base to have height at
least three: a hypothetical different example might have
several components permuted by the group, or nonlinear components.
The two lemmas above only sharpen the obstruction already found:
scalar linear centroid does not repair the abstract height-three
implication, while the same countermodel does not disprove a
stronger irreducibility theorem for the compatibility polynomial.
No numerical computation is required for either proof.
